% PM 트랙 Day 4 슬라이드 본문 — 손으로 쓴 정본(한 프레임 = 한 쪽). Day3_슬라이드_body.tex 와 같은 형식.
% 래퍼: Day4_슬라이드.tex(슬라이드판) / Day4_슬라이드_노트.tex(노트판, 우측에 \note).
% 화면에는 결론만, 설명·근거·예외는 각 프레임의 \note{}에.
% 그림은 ../그림/PM_Day4/*.pdf (벡터). 그림 번호 = 파일명 번호(그림 2-1 = fig_2_1_*). 모든 숫자는 Day4 워크북 완성 예시본(정본)에서.
% 데이터 일자 (a) 2027-02-05 런북·SAC 확정 → 실행 기록 02-28 / (b) 2027-05-28 G4. 컷오버 창 02-26 20:00~02-28 06:00(34h) · PoNR 01:40 · 롤백 최종 시한 16:00 · 트리거 4 발동 0 · 최종 AC 151.95 / VAC −5.13 / CPI 0.966.


\newcolumntype{L}[1]{>{\raggedright\arraybackslash}p{#1}}
\newcommand{\ra}{$\rightarrow$}
\newcommand{\til}{\textasciitilde}
\newcommand{\keymsg}[1]{\par\vfill{\color{mDarkTeal}\bfseries #1}\par}
\newcommand{\figcap}[1]{\par\smallskip{\footnotesize\color{black!60}#1}\par}
\newcommand{\fig}[2][0.66]{\centering\includegraphics[height=#1\textheight,width=\textwidth,keepaspectratio]{#2}\par}
\newcommand{\subhead}[1]{{\color{mDarkTeal}\bfseries #1}\par}
\newenvironment{tbl}[2][\small]{\begin{center}\usebeamercolor[fg]{normal text}\setlength{\tabcolsep}{4pt}\renewcommand{\arraystretch}{1.15}#1\begin{tabular}{#2}\toprule}{\bottomrule\end{tabular}\end{center}}

% =====================================================================
% 도입
% =====================================================================
\begin{frame}{PM 트랙 Day 4 — 오픈은 이벤트가 아니라 이관이다}
\begin{itemize}
\item 4일 특강 · 4일차 · 전환과 종료(Transition · Closing)
\item 사례는 같다 — 온담F\&B홀딩스 ONE ONDAM \textbf{P1}, BL-02 PMB 146.82 / 집행 한도 156.0
\item 시점이 둘 — 런북 확정 \textbf{2027-02-05} / 종료 \textbf{2027-05-28 G4}
\item 1차 컷오버 02-26 20:00\til{}02-28 06:00(34h), PoNR 01:40, 트리거 4 발동 0
\item 오늘의 문서는 프로젝트가 아닌 것을 향해 쓴다 — 결정 요청 칸이 아니라 \textbf{인수자 서명 칸}
\end{itemize}
\keymsg{오늘의 질문 — 오픈 다음 날 아침, 운영이 "우리는 인수한 적 없다"고 말하면 무엇이 없었던 것인가}
\note{표지. 교재 서문. Day3 워크북 완성 예시본(EX-01 대안 1 채택 10-16, MG3 02-19 회복)을 읽어 왔는지 확인한다. Day1은 결정하는 문서, Day2는 분해하고 숫자로 만드는 문서, Day3는 편차와 결정 요청의 문서였다면, Day4는 프로젝트 바깥의 사람이 읽고 서명하는 문서다 — 런북은 운영이 오픈 다음 날 읽고, 종료 보고서는 내부감사가 읽고, 편익 원장은 P1이 사라진 뒤 재무본부가 쓴다. 시점이 둘이라는 점(계획 / 실적 열 분리)을 워크북 §0.1 규칙 3으로 미리 못 박는다. 예상 소요 3분.}
\end{frame}

\begin{frame}{오늘의 흐름}
\begin{columns}[T]
\begin{column}{0.47\textwidth}
\subhead{오전 — 이론 (제1\til{}3장)}
\begin{itemize}
\item \textbf{1교시} 90분 · 전환·종료론과 컷오버 — 세 개의 시계 · 런북 · 게이트 · PoNR · 롤백 · 리허설 — 제1\til{}2장
\item \textbf{2교시} 90분 · 서비스 전환 — ITIL 판본 · SAC · ELS · SLA 3층 · error budget · 72시간 — 제3장
\item \textbf{3교시} 90분 · 종료 — 원가 3층 · 계약 종결 · 지체상금 0 · 이관 10건 — 제4장
\end{itemize}
\end{column}
\begin{column}{0.51\textwidth}
\subhead{오후 — 실습과 교훈·편익}
\begin{itemize}
\item \textbf{4교시} 90분 · \textbf{실습 ①} 컷오버 런북·SAC + 이관 체크리스트·SLA/OLA/UC — 워크북 §1\til{}2
\item \textbf{5교시} 75분 · 교훈과 편익 — Syllk · 여섯 칸 · ITO · 분모 정치 · 손익분기 66\% — 제5장
\item \textbf{6교시} 75분 · \textbf{실습 ②} 종료 보고서·교훈 등록부·편익 원장 — 워크북 §3\til{}5
\item \textbf{마무리} 30분 · 학술 배경 · 20종의 사슬을 거꾸로 · Product 트랙 예고 — 제6장
\end{itemize}
\end{column}
\end{columns}
\keymsg{오전은 "받는 쪽"의 문서를 읽고, 오후는 "받는 쪽"이 서명할 문서를 쓴다}
\note{교재 목차. 1교시가 두 장(제1·2장)을 묶는 이유는 컷오버가 종료론의 첫 실행이기 때문이다 — 런북의 마지막 행이 D+1 인계 서명이다. 2교시는 ITIL 어휘를 판본 중립으로 고정한 뒤 SAC·ELS·SLA를 다룬다. 3교시 종료 보고서는 "절감"과 "초과"가 동시에 참인 원가 3층이 핵심이다. 실습 ①은 런북 항목 3\til{}6(게이트·PoNR·트리거·검증)을 손으로 채우는 것이 중심이고, 실습 ②는 종료 보고서 항목 4·6과 편익 원장 측정 책임자 칸이다. 예상 소요 3분.}
\end{frame}

\begin{frame}{학습 목표}
\begin{columns}[T]
\begin{column}{0.49\textwidth}
\subhead{오전 — 읽을 수 있다}
\begin{itemize}
\item 프로젝트·운영·편익의 세 시계를 분리해 오픈의 위치를 말한다
\item 게이트를 "되돌리는 비용이 달라지는 곳"에 세우고 판정자를 배치한다
\item PoNR과 기술적 롤백 시한을 구별하고 결정권 상향 구간을 정한다
\item SAC·ELS 종료 조건·SLA 3층의 결손을 찾는다
\item 종료 보고서의 숫자를 성과보고·변경 로그의 행으로 추적한다
\end{itemize}
\end{column}
\begin{column}{0.49\textwidth}
\subhead{오후 — 쓸 수 있다}
\begin{itemize}
\item ⑯ 런북 게이트 3·PoNR·트리거 4·검증 6종의 조건을 관측값으로 적는다
\item ⑰ 이관 체크리스트의 인수 서명 4열과 SLA/OLA/UC 3층을 채운다
\item ⑱ 원가 3층(168/156/146.82) 대비 최종 AC 151.95를 분리해 적는다
\item ⑲ 교훈 6칸 중 "적용 대상·소유자" 칸을 실명으로 채운다
\item ⑳ 편익 원장의 측정 책임자·분모·손익분기를 서명란과 함께 만든다
\end{itemize}
\end{column}
\end{columns}
\keymsg{목표의 공통 척도 — 프로젝트가 사라진 뒤에도 그 문서가 누군가의 손에 남아 있는가}
\note{교재 서문의 "세 가지를 미리 말해 둔다"를 목표로 바꾼 것. 첫째, 오늘의 시나리오는 Day3 결정 요청에 스폰서가 답한 뒤의 세계다(EX-01 대안 1, R-19 0.60 배정, 대안 3 미발동). 둘째, 기준선 값을 바꾸지 말고 인용하라 — 종료에서도. 승인 예산 168 대비 16.05를 절감으로 적고 싶은 유혹이 가장 큰 시점이 종료다. 셋째, 모든 문서는 인수자 서명으로 끝난다 — 박준영 07:20, 네 사람의 인수 서명, 정미란 G4, 프로그램 PMO장 개정 요청 접수, 재경팀장 정본 인계. 예상 소요 3분.}
\end{frame}

\begin{frame}{Day3 통제 산출물 5종 \ra{} Day4 전환·종료 산출물 5종}
\begin{tbl}[\scriptsize]{L{3.0cm}L{5.4cm}L{2.3cm}}
Day4 산출물 & Day3에서 받는 입력 & 교재 \\
\midrule
⑯ 런북 · SAC & ⑫ EX-01 대안 1 · ⑭ 심각도 1·2 미해결 0 · ⑬ R-14(두 번 컷오버) · ⑪ CR-08 조건 ③·CR-11 & 제2장 · 3.2 \\
⑰ 이관 · SLA/OLA/UC · ELS & ⑬ 이슈 미결 6건 · ⑭ 인터페이스 47본 문서 · ⑮ 하도급 조건 ③·계약 변경 4·6·7번 · R-16 & 제3장 \\
⑱ 종료 보고서 · 이관 목록 · 계약 종결 & ⑫ BL-02(146.82)·EAC4 150.86 · ⑪ CR-01\til{}11 · ⑬ 잔여 EMV 7.26 · ⑮ 귀책 분해표 & 제4장 \\
⑲ 교훈 등록부 & ⑫ 조기 신호 5·착시 5 · ⑭ 감리 귀속 이의 · ⑬ IS-02 · AF-1-08 & 제5장 전반 \\
⑳ 편익 실현 원장 & ⑫ 편익 축 G(118억) · ⑬ R-12(두 분모) · ⑪ CR-04·CR-09 · Day1 ① BC §4·§6 & 제5장 후반 \\
\end{tbl}
\keymsg{20종의 사슬 — 런북이 시작이고 편익 원장이 끝이다. 되돌아가는 파선은 종결 기록이지 재기준선이 아니다}
\note{교재 서문의 표와 그림 0-1(마무리에서 그림 6-1로 다시 본다). 다섯 산출물은 서로를 인계한다 — 런북 D+1 판정 회의가 ⑰ 이관의 시작이고, SAC 조건부 1항목이 ⑱ 이관 목록 1번이 되며, 이관 6번(두 분모 판정, 나영선·2027-09 이사회)이 ⑳ 원장 BR-05로 간다. 강조할 것 — 종료 보고서의 모든 숫자는 Day3 문서의 어느 행에서 왔는지가 추적되어야 하고, 추적되지 않는 숫자는 종료 시점에 만들어진 숫자다(4.1절). 예상 소요 4분.}
\end{frame}

% =====================================================================
\section{1교시 (90분) — 전환·종료론과 컷오버}
% =====================================================================
\begin{frame}{이 교시에서 배우는 것}
\begin{itemize}
\item 프로젝트의 끝과 서비스의 시작은 같은 날이 아니다 — 세 개의 시계 (1.1)
\item 표준이 종료를 어떻게 정의하는가 — PRINCE2 5활동 · PMBOK 8 · ISO 21502 (1.2)
\item 임시조직의 끝은 제도로 정한다 — 해산은 사람별 일자로 (1.3)
\item 분야별 종료 게이트의 뼈대는 같다 — 증명서·하자 목록·보증·서명 (1.4)
\item 컷오버 — 런북·게이트 3·PoNR·롤백 트리거 4·검증 6종·리허설 3회 (제2장)
\end{itemize}
\keymsg{한 문장 — 오픈은 프로젝트가 시스템을 놓은 사건이 아니라, 운영이 받았다고 서명한 사건이어야 한다}
\note{교재 제1장·제2장 2.1\til{}2.7, 워크북 §1. 전반 40분은 종료론(제1장), 후반 50분은 컷오버(제2장)이며, 2.8 D$-$2 라이선스 사건은 2교시 첫머리에서 그림 2-6으로 다룬다. 이 교시의 실습 연결은 워크북 §1 항목 2\til{}6(시각표·게이트·PoNR·트리거·검증)이다. 수강생에게 미리 묻는다 — 최근 오픈에서 "운영이 받았다"는 서명은 누가 언제 했는가(제1장 생각해 볼 질문 1). 예상 소요 2분.}
\end{frame}

\begin{frame}{오픈 $\neq$ 종료 — 세 개의 시계}
\fig[0.58]{../그림/PM_Day4/fig_1_1_three_clocks.pdf}
\figcap{그림 1-1. 프로젝트 시계(G3 \ra{} 컷오버 \ra{} 2차 \ra{} ELS 종료 \ra{} G4 05-28)·운영 시계(ELS 착수 서명 07:20 \ra{} CAB \ra{} 하자보수 2028-05-28)·편익 시계(첫 측정 2027-04 \ra{} G5 2028-03-31). 오픈은 겹치는 점 하나다}
\keymsg{오픈은 세 시계가 만나는 점이 아니다 — 프로젝트의 마지막 마일스톤과 운영의 첫 달이 겹치는 곳이고, 편익 시계는 아직 시작하지도 않았다}
\note{교재 1.1. 온담의 프로젝트는 오픈(02-28 06:00) 뒤 석 달을 더 산다 — 2차 컷오버 03-27, ELS 종료 04-23, 운영 이관 착수 04-26, R5 05-07, 감리 G4 05-14, SAC 판정 05-21, 계약 종결 05-26, G4 05-28. 운영은 D+1 07:20 박준영의 ELS 착수 서명부터 시작되었고, 편익은 B-08 첫 측정(2027-04, 13.6시간)만 있고 판정은 G5다. "이관"의 정의 — 주는 쪽과 받는 쪽이 있고 받는 쪽이 받았다고 서명해야 끝난다. 런북 마지막 행이 "06:00 오픈"이 아니라 "D+1 07:20 서명"인 이유. 흔한 오류 — 헌장 종료일 = 컷오버일로 적어 PM의 마지막 석 달이 계약상 존재하지 않게 되는 것. 예상 소요 7분.}
\end{frame}

\begin{frame}{PRINCE2 7 Closing a Project — 5활동과 End Project Report}
\fig[0.56]{../그림/PM_Day4/fig_1_2_prince2_closing.pdf}
\figcap{그림 1-2. Prepare planned closure · Prepare premature closure · Hand over products · Evaluate the project · Recommend project closure — 4 산출물(End Project Report · Lessons Report · Benefits Management Approach · Follow-on Action Recommendations)이 ⑯\til{}⑳에 대응한다}
\keymsg{두 규율 — 인수는 운영의 공식 수용이지 PM의 선언이 아니다 / premature closure는 실패 처리가 아니라 정규 활동이다}
\note{교재 1.2. PRINCE2가 종료에서 주는 규율 두 가지를 강조한다. 첫째, hand over products의 인수(acceptance)는 운영·유지보수 조직의 공식 수용 — ⑯⑰의 서명이 그것이다. 둘째, premature closure는 정규 활동이며 중단에도 인계·평가·권고가 있다. P1 헌장 §14의 중단 조건(편익 전망 손익분기 66\% 미만, 컷오버 창 2회 연속 이탈)과 손실 확정 절차(재경팀장 30일 내)는 이 활동의 사전 문서화이고, 종료 보고서에는 "미발동"으로 기록이 남는다. 흔한 오류 — End Project Report 목차를 그대로 옮기고 절마다 한 문단씩 채우는 것. 표준은 무엇이 있어야 하는가를 말하지 그 값이 어디서 오는가를 말하지 않는다. 예상 소요 7분.}
\end{frame}

\begin{frame}{PMBOK 8판 Closing Focus Area · ISO 21502 §7}
\begin{itemize}
\item PMBOK 8 Closing — 인도물 이전·계약 종결·재무 정산·팀 해산·기록·성과 평가·교훈·\textbf{편익 계획 이전}
\item 6판 4.7 한 프로세스 \ra{} 8판은 계약 종결·재무 정산을 Closing 안에 명시
\item 가장 무거운 문장 — "편익 실현 계획을 프로젝트 후 조직에 이전한다" \ra{} ⑳ 소유 재무본부
\item 운영 준비도(operational readiness)를 종료 조건으로 적는다 \ra{} SAC·⑰
\item ISO 21502 §7 — 종료 조건 확인·인계·계약 종결·기록·교훈·편익 책임 명시. 계약이 인용할 때 검수 조건의 해석 기준
\end{itemize}
\keymsg{세 표준의 공통분모 — "누구에게 넘기는가"가 적혀야 종료다. 편익 계획도, 미해결 사항도, 잔여 리스크도}
\note{교재 1.2. PMBOK 6판의 편익 관리 계획(1.2.6.2)에는 "누구에게 넘기는가"가 없었고 8판은 넘긴다고 쓴다 — 온담 편익 원장의 "소유 재무본부(재경팀장), P1 종료 후 정본 관리자" 한 줄이 그 실행이다. ISO 21502가 실무에서 중요한 순간은 계약서가 "ISO 21502에 따른 프로젝트 관리"를 인용할 때이며, 그때 §7이 검수 조건의 해석 기준이 된다 — 원문 대조를 전제로 말한다(판본은 검색으로 확인). 2026년은 PMBOK 8·ITIL v5·개인정보보호법 개정의 3중 전환기이며, 판본 이름을 외우지 말고 어휘를 판본 중립으로 쓰라는 것이 2교시 3.1의 태도다. 예상 소요 6분.}
\end{frame}

\begin{frame}{분야별 종료 게이트 — 이름은 다르고 뼈대는 같다}
\begin{tbl}[\scriptsize]{L{1.5cm}L{3.6cm}L{2.2cm}L{1.9cm}L{1.6cm}}
분야 & 단계적 증명서 & PoNR / go 조건 & 하자·보증 & 받는 쪽 \\
\midrule
건설·EPC & MC \ra{} 시운전 \ra{} 성능 시험 \ra{} PAC \ra{} FAC & 성능 시험 통과 & punch list A/B/C · DLP & 발주자 \\
SI & 릴리스 검수 \ra{} 최종 검수 \ra{} 잔금·유보금 & 데이터 정본 전환 & 하자보수 1년 · 유지보수 별도 & 운영 조직 \\
제조 NPD & 파일럿 런 \ra{} Run@Rate \ra{} 양산 이관 & 수율·불량률 & ECN backlog & 생산 \\
제약·임상 & 데이터 실사 \ra{} 쿼리 해소 \ra{} DB lock & \textbf{DB lock(규제가 명시)} & — & 규제 제출 \\
\end{tbl}
\keymsg{어휘가 아니라 구조를 빌린다 — 조건부 인수의 정규화(EPC) · PoNR의 명시적 정의(제약) · go 조건의 숫자화(NPD)}
\note{교재 1.4. 뼈대는 다섯 — 단계적 증명서, 하자 목록, 성능 확인, 보증 기간, 받는 쪽의 서명. SAC에 해당하는 것은 EPC의 PAC 조건 목록이고, 미해결 사항 이관 목록은 punch list B(인수 후 기한 내 해소)다. SI가 EPC와 다른 점은 컷오버가 데이터·인터페이스·단말의 세 갈래로 동시에 일어난다는 것이며, 지체상금은 기준선 완료일 대비 지연 일수이므로 Day3의 귀책 분해표가 매월 누적되어야 종료 보고서 한 줄이 된다(3교시 4.3). 제약의 교훈 — 되돌릴 수 없는 시각을 규제가 정의한다. SI의 PoNR은 대개 암묵적이고, 그것이 2.3절의 주제다. 흔한 오류 — SI 계약에 FAC, EPC 계약에 하자보수를 쓰는 것. 예상 소요 6분.}
\end{frame}

\begin{frame}{임시조직의 끝 — Day1의 Lundin \& Söderholm을 회수한다}
\begin{itemize}
\item 전환(Transition) — 이후 상태는 영구조직의 것 \ra{} 전환 완료의 판정권도 영구조직(SAC = 박준영)
\item 제도화된 종료 — 종료 보고서·해산·복귀가 제도로 정해져 있다
\item 두 실무 현상 — 종료 직전 인력 이탈(IS-04·IS-11) · 교훈의 실패(조직이 해산하면 배운 것도 해산)
\item Staw의 처방 — 끝났다고 선언하는 사람과 받았다고 서명하는 사람은 달라야 한다
\item 해산은 사람별 일자로 — 발주사 14명 중 12명 06-01 복귀·2명 P5 / 수행사 05-31 철수·하자보수 4명 비상주
\end{itemize}
\keymsg{"05-31 해산" 한 줄이면 6월 첫 주의 하자 문의는 아무에게도 가지 않는다}
\note{교재 1.3. Day1 1.3절의 네 구조 개념(시간·과업·팀·전환)·네 행위 개념 중 전환과 제도화된 종료를 회수한다. Staw(1976)는 Day3 6.6에서 계속 투자의 처방으로 빌렸다 — 중단의 종료에서는 착수 결정자(김선호·프로그램 이사회)가 아닌 판정자(재경팀장 손실 확정, 프리모템 F3)가, 정상 종료에서는 착수 결정자가 "끝났다"를 선언하고 싶어 하므로 받는 쪽(박준영)의 서명이 판정이 된다. 온담의 OD-POS 2.1 종료 세션(03-19, 박세연 CIO 주관, 15년 기록 아카이브)은 감상이 아니라 Bridges의 "끝냄" 처리이며 2교시 3.9에서 다시 본다. P1 PMO 3명 중 1명 하자보수기 잔류는 "승인 전"으로 이관 목록에 남는다. 예상 소요 6분.}
\end{frame}

\begin{frame}[plain]{}
\fig[0.86]{../그림/PM_Day4/fig_2_1_runbook_timeline.pdf}
\figcap{그림 2-1. 1차 컷오버 런북 시각표 T$-$24h\til{}D+1 — 네 갈래(C·D·I·S) 38행, 게이트 1·2·3 계획/실제(+5·+38·+30), PoNR 01:40, 롤백 최종 시한 16:00, 오픈 판정 04:30 \ra{} 오픈 06:00 \ra{} D+1 07:20 서명. 여유 2h 20m 중 1h 10m 소진}
\keymsg{런북은 시각(時刻)이 주어인 유일한 프로젝트 문서 — 확인 열은 "정상 확인"이 아니라 관측값(건수·해시·성공률)으로}
\note{교재 2.1, 워크북 §1 항목 2. T0 = 02-26(금) 20:00 거래 정지, 창 34h(P-35, 허용범위 0). T$-$24h가 있는 이유 — 전날의 최종 상태 점검·전 백업·컨트롤룸 개설·외부 통보(여덟 행)가 담당자의 체크리스트에만 있으면 그 담당자의 병가로 창 시작 전에 이미 늦는다. 행 7(임시 라이선스 유효 확인)은 D$-$2 사건이 없었다면 없었을 행이다. 갈래를 나누는 것은 병렬 작업을 시각표 위에서 보이게 하고 부분 롤백(RB-3)을 정의하기 위해서다. 실행 기록 열 — 편차 총계 게이트 2 +38·게이트 3 +30·이력 적재 $-$25·배포 +10·오픈 0, 여유 소진 1h 10m. 이 숫자가 2차(03-27)·R5·P2 런북이 이번 런북을 베낄 수 있게 한다. 전체 폭 한 장이므로 상단·하단을 순서대로 짚는다. 예상 소요 8분.}
\end{frame}

\begin{frame}{go/no-go 3게이트와 PoNR 01:40 — 되돌리는 비용이 달라지는 곳}
\begin{tbl}[\scriptsize]{L{1.9cm}L{1.9cm}L{3.9cm}L{1.6cm}L{1.6cm}}
게이트 & 계획 / 실제 & 조건(전부 충족) & 판정자 & 되돌리는 비용 \\
\midrule
1 착수 & 19:30 / 19:35 & 인력 97명·백업 해시·임시 라이선스 유효·CAB 잠금 & P1 PM(위임) & 0 \\
2 정본 전환 & 01:00 / \textbf{01:38} & 마스터 검증 6종 통과·재실행 2회 이내 & \textbf{박세연} & 시간(재기동 2h) \\
\textbf{PoNR} & 01:40 / 01:40 & 신 원장 '정본' 활성화 & — & 시간 + 손실 \\
3 매장 배포 & 12:00 / \textbf{12:30} & 이력 6종·IF 47/47·스모크·RB-2 미발동 & \textbf{정미란} & + 매장 재방문 2.1억 \\
기술적 시한 & 16:00 & 오픈 06:00 $-$ 14h(복원 8h+검증 2h+캐시 4h) & 정미란 & 불가 \\
\end{tbl}
\keymsg{세 게이트를 모두 스폰서에게 맡기면 새벽 1시의 게이트 2는 "스폰서 깨우기 vs 그냥 진행"이 된다 — 데이터는 CIO, 돈과 매장은 스폰서}
\note{교재 2.2\til{}2.3, 워크북 §1 항목 3·4. 게이트의 세 요소 — 조건·판정자·시간 한도. 게이트 2의 38분이 핵심 — 00:30 검증 1차에서 표본 2,000행 중 1건 불일치(가격 마스터 VAT 소수 자릿수). 통과 기준 0과 재실행 허용 2회가 미리 적혀 있었기 때문에 새벽 1시에 토론하지 않고 재실행했다. 게이트 3의 30분은 스폰서의 열람 시간 — "스폰서가 늦었다"가 아니라 "판정자 열람 30분을 계획에 포함"으로 적어야 다음 런북이 배운다. PoNR은 기술적으로 되돌릴 수 없는 시각이 아니라 되돌리면 잃는 것이 생기는 시각 — 01:40부터 이력 적재·03:00 3PL 폴백 배치 출고 지시·매장 캐시가 신 원장에 기록된다. 01:40\til{}16:00의 14h 20m이 "되돌릴 수는 있으나 잃는 구간"이고 롤백 결정권이 PM에서 정미란(부재 시 김선호)으로 한 층 올라간다. 게이트 2 직후에 PoNR을 둔 이유 — 떼어 놓으면 03:00 폴백 배치가 구 원장 기준으로 나가 이중 원장이 된다. 예상 소요 9분.}
\end{frame}

\begin{frame}{롤백 트리거 4 · 롤백의 비대칭 — 롤백도 프로젝트다}
\fig[0.50]{../그림/PM_Day4/fig_2_2_reversal_cost.pdf}
\figcap{그림 2-2. 되돌리는 비용의 계단 — 게이트 1 전 0 \ra{} 시간만 \ra{} 시간+손실(PoNR) \ra{} +매장 방문 2.1억(게이트 3) \ra{} 불가(16:00). 트리거 RB-1\til{}4의 위치와 관측값(1/2 · 08:35/11:00 · 9/31점 · 100\%/99.0\%) — 발동 0}
\keymsg{임계값이 숫자가 아니면 판정은 토론이 된다 — 발동하지 않은 트리거의 관측값이 다음 컷오버의 임계값 근거다}
\note{교재 2.4, 워크북 §1 항목 5(그림 2-3이 트리거 비율 막대). 트리거의 요소 — 지표·임계값·관측 시각(런북 어느 행)·판정자·발동 시 조치·복구 소요. RB-1 마스터 검증 미통과(재실행 2회 후, 박세연, 전체 롤백 손실 0, 다음 창 대안 2 03-12) / RB-2 이력 적재 11:00 미완(PM \ra{} PoNR 후 정미란, 손실 있는 롤백 vs 오픈 연기 10:00) / RB-3 배포 실패 31점(5\%) 또는 직영 10점(오정근·박세연, 매장 갈래 부분 롤백) / RB-4 스모크 99.0\% 미만·PG 1본 실패(정미란, 오픈 연기 \ra{} 부분 운영). RB-2 임계값 11:00은 리허설 환산 12:30보다 보수적 — "언제 실패인가"가 아니라 "언제까지 결정하면 나머지 계획이 살아 있는가". RB-3 판정자가 오정근인 이유 — 매장 재방문 6영업일·2.1억은 영업본부의 비용이므로 Senior User가 판정한다. 비대칭 — 컷오버 계획은 38행, 롤백 계획은 대개 한 줄. PoNR 이후 롤백은 신 원장 거래의 역대사·출고 지시 취소·펌웨어 회수가 있어 컷오버보다 어렵다. 예상 소요 8분.}
\end{frame}

\begin{frame}{데이터 검증 6종 — 이 데이터의 정답을 누가 판정하는가}
\begin{tbl}[\scriptsize]{L{2.2cm}L{2.4cm}L{2.0cm}L{2.0cm}L{2.3cm}}
검증 & 분모 & 리허설 1 (11-23) & 본 컷오버 & 서명 \\
\midrule
① 건수 & 1,842 테이블 & 0 & 0 / 0 & 용역사 PM·PMO 품질 \\
② 금액 합계 & 5년 원장 & $-$41,200원 \ra{} 규칙 정정 & 0 / 0 & \textbf{재경팀장} \\
③ 키 유일성 & PK 1,842 & \textbf{1,204}(브랜드 간 중복, R-10) & 0 / 0 & \textbf{박세연} \\
④ 참조 정합성 & FK 3,916 & \textbf{412}(폐점·단종) & 0 / 0 & 박세연 \\
⑤ 코드 매핑 & 4,318 코드 & \textbf{96}(97.8\%) & 0 / 0 & 수행 통합팀·PMO \\
⑥ 표본 대조 & 2,000 / 5,000행 & 18 & \textbf{1 \ra{} 재실행 \ra{} 0} / 0 & PMO 품질·용역사 PM \\
\end{tbl}
\keymsg{분모 · 통과 기준 0 · 서명자 — 셋 중 하나가 없으면 전환 검증은 자기 채점이다. 리허설 1회차의 미통과 3종은 정본 판정(G2)의 입력이었다}
\note{교재 2.5, 워크북 §1 항목 6. 전환 대상 1,842테이블 7.4TB \ra{} CR-11(5년 필터) 5.1TB, 마스터 6종 2,184,310행·이력 3.21억 행. 여섯 각도 — 건수·금액 합계·키 유일성·참조 정합성·코드 매핑 커버리지·표본 필드 대조. 통과 기준이 "1\% 이내"이면 그 1\%가 어느 테이블에 몰려 있는지를 묻지 않게 된다. 리허설 1회차의 1,204·412·96은 전환 도구의 문제가 아니라 원본 데이터의 정본이 정해지지 않은 문제였고, G2(11-27) 마스터 6종 정본 판정으로 사라졌다 — Day3의 데이터 정본 판정이 여기서 회수된다. 본 컷오버 표본 1건은 성격이 다르다(리허설 3회차 통과 후 발생) — 리허설 시점의 통과를 컷오버 시점의 통과로 알지 말 것. 서명 열 — 재경팀장이 금액 합계에 서명하지 않으면 첫 월마감의 "합계가 다르다"가 전환 오류인지 운영 오류인지 가릴 수 없다. 예상 소요 8분.}
\end{frame}

\begin{frame}{리허설 3회 · 두 번 컷오버 — 결과가 나빠야 좋은 리허설}
\fig[0.52]{../그림/PM_Day4/fig_2_4_rehearsals.pdf}
\figcap{그림 2-4. 1회차 11-23 41h 20m(7.4TB 전체, 창 초과 7h 20m, 미통과 3종) \ra{} 2회차 01-08 33h 10m(5.1TB, 표본 3건) \ra{} 3회차 01-22 31h 40m(롤백 리허설 7h 50m 병행) \ra{} 본 컷오버 32h 50m 상당 / 창 34h}
\keymsg{회차마다 목적이 다르다 — 1회 무엇이 깨지는가 · 2회 창 안에 들어오는가 · 3회 여유와 PoNR 이후 롤백의 실측. 리허설 없이 정한 PoNR은 추측이다}
\note{교재 2.6\til{}2.7, 워크북 §1 항목 7. 리허설의 목적은 통과가 아니라 실측 — 소요 시간, 실패하는 검증, 병목 갈래, 교대 인수 누락. 1회차를 전체 볼륨으로 돌린 덕에 "필터 없이는 창에 들어오지 않는다"가 실측되어 CR-11의 근거가 되었다. 축소 볼륨 리허설의 함정 — 인덱스·락·I/O는 비선형이라 10\%로 3시간이면 전체 30시간이 아니다. 3회차의 롤백 병행(7h 50m, 대사 1,140건)이 기술적 시한 16:00과 결정권 상향의 근거다. 두 번 컷오버(R-14) — 1차 매장·주문 34h / 2차 물류 03-27 22:00\til{}06:00 8h, 나눈 기준은 조직이 아니라 P2 설비 준비 상태(리드타임 38주). 비용 셋 — 임시 상태 4주(매장 = 신 원장, 센터 = 구 SAP, 폴백 배치 2회)·"하지 말 것" 2줄·2차 게이트 1 조건 4개. 1차 하이퍼케어 심각도 2 이상 ≤ 5(실제 4)가 2차 게이트 1 조건 ①이었던 것이 L-15 — 두 컷오버를 달력이 아니라 안정화 지표로 묶는다(그림 2-5). 예상 소요 7분.}
\end{frame}

\begin{frame}{1교시 핵심 정리}
\begin{itemize}
\item 세 개의 시계 — 오픈은 프로젝트의 마지막 마일스톤과 운영의 첫 달이 겹치는 점, 편익 시계는 아직 전
\item 이관 = 받는 쪽의 서명 — PRINCE2 hand over · PMBOK 8 편익 계획 이전 · 해산은 사람별 일자로
\item 게이트는 되돌리는 비용이 달라지는 곳에 — 조건·판정자·시간 한도, 데이터는 CIO·돈은 스폰서
\item PoNR 01:40 $\neq$ 기술적 시한 16:00 — 그 사이 14h 20m의 롤백 결정권은 한 층 위
\item 트리거·검증은 숫자와 서명자로 — 발동 0의 관측값과 리허설 1회차의 실패가 다음 컷오버의 자산
\end{itemize}
\keymsg{컷오버는 되돌릴 수 없는 시각을 정하는 기술이고, 종료는 받는 쪽이 서명하는 문서를 쓰는 기술이다}
\note{교재 제1장·제2장 핵심 정리. 세 오류를 세 문서의 부재로 다시 말한다 — 오픈 = 종료(SAC 부재), 보고서 없는 해산(증명 부재, 2019 스마트오더 세 줄), 소멸하는 미해결 사항(이관 목록 부재, 소유자 실명·기한·승인 여부). 종료에서 양식의 부재는 프로젝트 안에서 발견되지 않는다 — 프로젝트가 사라진 뒤 운영·내부감사·다음 PM이 발견한다. 2교시로 넘어가며 — D$-$2 라이선스 사건(2.8, 그림 2-6)이 게이트 1 조건에 행 하나를 추가하게 만든 인과 사슬과, ITIL 어휘를 판본 중립으로 고정하는 이유(3.1)부터 시작한다. 휴식 전 워크북 §1 항목 3·4·5의 빈 템플릿을 펴 두라고 안내한다. 예상 소요 3분.}
\end{frame}
